\documentclass[aspectratio=169,10pt]{beamer}
\usetheme{Madrid}
\usecolortheme{default}

\usepackage{hyperref}
\usepackage{tikz}
\usepackage{amsmath,amssymb}

% --- Custom Colors ---
\definecolor{unibblue}{RGB}{0,51,102}
\definecolor{unibgold}{RGB}{212,175,55}
\definecolor{unibgreen}{RGB}{0,128,0}

\setbeamercolor{palette primary}{bg=unibblue,fg=white}
\setbeamercolor{palette secondary}{bg=unibblue!85,fg=white}
\setbeamercolor{palette tertiary}{bg=unibblue!70,fg=white}
\setbeamercolor{structure}{fg=unibblue}
\setbeamercolor{block title}{bg=unibblue!85,fg=white}
\setbeamercolor{block body}{bg=unibblue!10,fg=black}
\setbeamercolor{alertblock title}{bg=unibgold!90,fg=black}
\setbeamercolor{alertblock body}{bg=unibgold!15,fg=black}

\title{Persamaan Diferensial}
\subtitle{Minggu 12: Persamaan Panas 1D \& Manajemen Termal Konduktor Daya}
\author{Novalio Daratha \& Muhammad Arfan}
\institute{Program Studi Teknik Elektro\\Universitas Bengkulu}
\date[\today]{2026 \\ \vspace{2mm}\small\textit{Tanggal Kompilasi: \today}}

\begin{document}

\begin{frame}
  \titlepage
\end{frame}

\begin{frame}{Capaian Pembelajaran (CPMK 4)}
  \begin{itemize}
    \item \textbf{CPMK 4.1:} Memahami penurunan Hukum Konduksi Fourier pada konduktor kabel berarus.
    \item \textbf{CPMK 4.2:} Menyelesaikan PDP Panas 1D untuk kondisi batas homogen (*Dirichlet cooling*) dan terisolasi (*Neumann*).
    \item \textbf{CPMK 4.3:} Menganalisis dekomposisi suhu transien dan distribusi suhu keadaan mantap (*steady-state*).
    \item \textbf{CPMK 4.4:} Menerapkan skema beda hingga FTCS pada simulasi Julia (Simulasi Persamaan Panas 1D).
  \end{itemize}
\end{frame}

\begin{frame}{Daftar Isi}
  \tableofcontents
\end{frame}

\section{Pemodelan Termal Konduktor Listrik}
\begin{frame}{Penurunan Persamaan Difusi Panas 1D}
  Berdasarkan Hukum Konduksi Kalor Fourier ($q = -k \frac{\partial u}{\partial x}$) dan konservasi energi pada batang konduktor dengan difusivitas termal $\alpha = \frac{k}{\rho c_p}$:
  \begin{alertblock}{Persamaan Panas 1D}
    $$ {\color{unibblue} \frac{\partial u}{\partial t} = \alpha \frac{\partial^2 u}{\partial x^2}} $$
    \begin{itemize}
      \item $u(x,t)$: Distribusi temperatur pada posisi $x$ dan waktu $t$ ($^\circ\text{C}$).
      \item $\alpha$: Difusivitas termal material ($m^2/s$).
    \end{itemize}
  \end{alertblock}
  \pause
  \textit{Aplikasi Elektro: Perhitungan ampacity (kapasitas hantar arus) maksimum kabel bawah tanah dan disipasi panas transformator.}
\end{frame}

\section{Solusi Analitik Deret Fourier}
\begin{frame}{Penyelesaian dengan Syarat Batas Dirichlet}
  Untuk kabel panjang $L$ dengan kedua ujung dijaga pada suhu lingkungan $0^\circ\text{C}$ ($u(0,t)=0, u(L,t)=0$) dan profil suhu awal $u(x,0) = f(x)$:
  \pause
  \begin{block}{Solusi Lengkap}
    $$ {\color{unibblue} u(x,t) = \sum_{n=1}^{\infty} B_n \sin\left(\frac{n\pi x}{L}\right) e^{-\lambda_n^2 \alpha t}}, \quad \lambda_n = \frac{n\pi}{L} $$
  \end{block}
  \pause
  \textbf{Faktor Peluruhan Eksponensial:}
  $$ \tau_n = \frac{1}{\alpha \lambda_n^2} = \frac{L^2}{n^2 \pi^2 \alpha} $$
  \textit{Harmonik spasial frekuensi tinggi ($n=2,3,\dots$) meluruh jauh lebih cepat, menyisakan gelombang sinusoidal mulus mode fundamental ($n=1$).}
\end{frame}

\section{Kondisi Batas Non-Homogen \& Keadaan Mantap}
\begin{frame}{Transien vs Keadaan Mantap (*Steady-State*)}
  Jika kedua ujung kabel memiliki suhu konstan berbeda: $u(0,t) = T_1$ dan $u(L,t) = T_2$:
  \pause
  Solusi total dipecah menjadi:
  $$ u(x,t) = u_{\text{steady}}(x) + u_{\text{transient}}(x,t) $$
  \pause
  \begin{block}{Distribusi Suhu Keadaan Mantap (Saat $t \to \infty$)}
    $$ \frac{d^2 u_{\text{steady}}}{dx^2} = 0 \implies {\color{unibblue} u_{\text{steady}}(x) = T_1 + \frac{T_2 - T_1}{L} x} \quad \text{(Garis Lurus Linear)} $$
  \end{block}
\end{frame}

\section{Kuis \& Simulasi Komputasi}
\begin{frame}{Simulasi Numerik Julia: FTCS}
  \begin{block}{Simulasi Persamaan Panas 1D}
    $$ u_i^{n+1} = u_i^n + r (u_{i+1}^n - 2u_i^n + u_{i-1}^n), \quad r = \frac{\alpha \Delta t}{\Delta x^2} \le 0.5 $$
    \begin{itemize}
      \item Mempelajari syarat stabilitas von Neumann ($r \le 0.5$).
      \item Visualisasi difusi titik panas (*hot-spot*) kabel daya.
    \end{itemize}
  \end{block}
\end{frame}

\begin{frame}{Kuis Interaktif 12}
  \textbf{Soal:} Jika panjang kabel tembaga diduplikasi menjadi $2\times$ lipat semula ($L' = 2L$), berapakah peningkatan waktu pendinginan $\tau$ yang dibutuhkan kabel hingga mencapai suhu dingin?
  \pause
  \begin{block}{Jawaban:}
    Karena konstanta waktu peluruhan $\tau \propto L^2$:
    $$ \tau' \propto (2L)^2 = 4 L^2 = 4\tau $$
    Waktu pendinginan meningkat \textbf{4 kali lipat}!
  \end{block}
\end{frame}

\begin{frame}{Rangkuman}
  \begin{enumerate}
    \item Persamaan panas adalah PDP tipe parabolik yang memodelkan proses disipasi irreversibel.
    \item Solusi selalu konvergen menuju gradien temperatur linear pada keadaan mantap.
    \item Kestabilan numerik FTCS menuntut langkah waktu $\Delta t \le \frac{\Delta x^2}{2\alpha}$.
  \end{enumerate}
\end{frame}

\end{document}
